\section{Coefficient compatibility and feasibility}
\label{sec:compatibility}

\subsection{From coefficient kernels to a matrix-space condition}

Set
\begin{align}
	W_0   & =P_{00},                                     &
	W_1   & =P_{01},                                     &
	W_2   & =P_{10},                                     &
	W_3   & =P_{11},                                     &
	W_4   & =T,
	\nonumber                                              \\
	\mc W & =\operatorname{span}\{W_0,W_1,W_2,W_3,W_4\}.
	\label{eq:W}
\end{align}
The corresponding physical basis is
\begin{align}
	\mathsf B_0 & =J_1(E_{00}\oplus(-1)),
	            &
	\mathsf B_1 & =J_1(E_{01}\oplus0),
	\nonumber                             \\
	\mathsf B_2 & =J_1(E_{10}\oplus0),
	            &
	\mathsf B_3 & =J_1(E_{11}\oplus(-1)),
	\nonumber                             \\
	\mathsf B_4 & =J_1(0\oplus1).
	\label{eq:B}
\end{align}
Using \cref{eq:Q}, the one-site and two-site letters become
\begin{align}
	M^{\alpha\beta}
	 & =\sum_{a=0}^4(W_a)_{\alpha\beta}\mathsf B_a,
	\nonumber                                                                 \\
	G^{\alpha\beta}
	 & =\sum_{\gamma,a,b}
	    (W_a)_{\alpha\gamma}(W_b)_{\gamma\beta}\mathsf B_a\otimes \mathsf B_b
	\nonumber                                                                 \\
	 & =\sum_{a,b=0}^4(W_aW_b)_{\alpha\beta}\mathsf B_a\otimes \mathsf B_b.
	\label{eq:blocked-coordinates}
\end{align}
The $\mathsf B_a$ are independent, as are their 25 tensor products.

For complex coefficients $c_{\alpha\beta}$, define
\begin{align}
	F_1(c) & =\sum_{\alpha,\beta}c_{\alpha\beta}M^{\alpha\beta},
	       &
	F_2(c) & =\sum_{\alpha,\beta}c_{\alpha\beta}G^{\alpha\beta}.
	\label{eq:coefficient-maps}
\end{align}
Write
$\mc W^2=\operatorname{span}\{W_aW_b:0\leq a,b\leq4\}$.
The kernels in \cref{eq:coefficient-maps} are the annihilators of
$\mc W$ and $\mc W^2$ under the nondegenerate bilinear entry pairing.
Taking annihilators reverses inclusions, so
\begin{align}
	\ker F_1\subseteq\ker F_2
	 & \quad\Longleftrightarrow\quad
	\mc W^2\subseteq\mc W,
	\nonumber                        \\
	\ker F_1=\ker F_2
	 & \quad\Longleftrightarrow\quad
	\mc W^2=\mc W.
	\label{eq:kernel-criterion}
\end{align}
A linear refinement map on the letter span exists exactly when
$\ker F_1\subseteq\ker F_2$: the assignment
$F_1(c)\mapsto F_2(c)$ is well defined precisely under this inclusion.
Thus closure of $\mc W$ under multiplication is necessary even
before asking for an inverse map.

For reference, these tests have fully explicit reduced coordinates.
Let $h_{j,uv}=(H_j)_{uv}$ for $j=0,1$, and let
$(j_a,k_a)$ run through $(0,0),(0,1),(1,0),(1,1)$ for $a=0,1,2,3$.
Then
\begin{align}
	w_{a;(u,u'),(v,v')}
	 & :=(W_a)_{(u,u'),(v,v')}
	=\sqrt{p_up_{u'}}
	h_{j_a,uv}\overline{h_{k_a,u'v'}},
	\nonumber                        \\
	w_{4;(u,u'),(v,v')}
	 & =p_u\delta_{uu'}\delta_{vv'}.
	\label{eq:entry-coordinates}
\end{align}
Every product coefficient is obtained by the four-term contraction
\begin{align}
	z_{ab;\alpha\beta}
	 & :=(W_aW_b)_{\alpha\beta}
	\nonumber                              \\
	 & =w_{a;\alpha,00}w_{b;00,\beta}
	+w_{a;\alpha,01}w_{b;01,\beta}
	\nonumber                              \\
	 & \quad+w_{a;\alpha,10}w_{b;10,\beta}
	+w_{a;\alpha,11}w_{b;11,\beta}.
	\label{eq:product-entries}
\end{align}
Closure would require each of these matrices to be a linear
combination of the five $W_c$.

The putative coefficients are determined, not independent.
Define
\begin{align}
	\Gamma_{cd}
	 & =\sum_{\alpha,\beta}
	\overline{w_{c;\alpha\beta}}w_{d;\alpha\beta},
	 &
	b_{d;ab}
	 & =\sum_{\alpha,\beta}
	\overline{w_{d;\alpha\beta}}z_{ab;\alpha\beta}.
\end{align}
The Gram matrix $\Gamma$ is positive definite. Taking inner products
of $W_aW_b=\sum_c\mu_{ab}^cW_c$ with $W_d$ gives
$\sum_c\Gamma_{dc}\mu_{ab}^c=b_{d;ab}$, hence
\begin{align}
	\mu_{ab}^c
	 & =\sum_{d=0}^4(\Gamma^{-1})_{cd}b_{d;ab},
	\nonumber                                   \\
	z_{ab;\alpha\beta}
	 & \stackrel{\mathrm{closure}}{=}
	\sum_{c=0}^4\mu_{ab}^cw_{c;\alpha\beta}.
	\label{eq:closure-equations}
\end{align}
The following argument proves that these closure equations have no
solution anywhere in \cref{eq:parameter-space}.

\subsection{The rank-one transfer forces two-dimensional orbits}

In the normalization \cref{eq:normal-form}, put
\begin{align}
	r                           & =\operatorname{vec}(R),
	                            &
	\ell\,\operatorname{vec}(Y) & =\tr Y.
\end{align}
Then $T=r\ell$ and $\ell r=1$.
Define its left column orbit and right row orbit by
\begin{align}
	\mc U_{\mathrm{col}} & =\mc W r,
	                     &
	\mc U_{\mathrm{row}} & =\ell\mc W,
	                     &
	a                    & =\dim\mc U_{\mathrm{col}},
	                     &
	b                    & =\dim\mc U_{\mathrm{row}}.
	\label{eq:orbits}
\end{align}
Under vectorization and trace duality, these are
\begin{align}
	\mc U_{\mathrm{col}}
	 & \simeq\operatorname{span}
	\{R,A_jRA_k^\dagger:j,k=0,1\},
	\nonumber                    \\
	\mc U_{\mathrm{row}}
	 & \simeq\operatorname{span}
	\{\Id_2,A_k^\dagger A_j:j,k=0,1\}.
	\label{eq:orbit-matrices}
\end{align}
The second identification follows from
$\tr(A_jYA_k^\dagger)=\tr(A_k^\dagger A_jY)$.

\begin{lemma}
	\label{lem:orbit}
	If $\mc W$ is closed under multiplication, then $a=b=2$.
\end{lemma}

\begin{proof}
	First, both dimensions are at least two. If $a=1$, all matrices
	$A_jRA_k^\dagger$ are proportional to $R$.
	Since $A_0\neq0$ and $R>0$, the identity
	$A_0RA_0^\dagger=cR$ has $c>0$ and makes $A_0$ invertible.
	Writing $A_0RA_1^\dagger=dR$ gives
	\begin{align}
		A_0^\dagger & =cR^{-1}A_0^{-1}R,
		            &
		A_1^\dagger & =dR^{-1}A_0^{-1}R,
		\nonumber                        \\
		A_1^\dagger & =(d/c)A_0^\dagger,
	\end{align}
	contrary to independence of $A_0,A_1$.

	If $b=1$, there are scalars $c>0$ and $d$ such that
	$A_0^\dagger A_0=c\Id_2$ and $A_0^\dagger A_1=d\Id_2$.
	They imply
	\begin{align}
		A_0 & =c(A_0^\dagger)^{-1},
		    &
		A_1 & =d(A_0^\dagger)^{-1}
		=(d/c)A_0,
	\end{align}
	which is the same contradiction.

	Now assume closure. Since $T\in\mc W$, the space
	\begin{align}
		\mc I
		 & :=\operatorname{span}\{XTY:X,Y\in\mc W\}
		\nonumber                                   \\
		 & =\operatorname{span}
		\{(Xr)(\ell Y):X,Y\in\mc W\}
		\label{eq:ideal}
	\end{align}
	is contained in $\mc W$.
	Outer products of independent column and row bases are independent:
	applying a dual functional to the columns first leaves a linear
	combination of the independent rows, whose coefficients must then
	vanish. Therefore
	\begin{align}
		ab=\dim\mc I\leq\dim\mc W=5.
	\end{align}
	The integers $a,b$ are both at least two, so $a=b=2$.
\end{proof}

Under the closure hypothesis, $\mc I$ would be a four-dimensional
two-sided ideal. Indeed, for $Z,X,Y\in\mc W$,
\begin{align}
	Z(XTY) & =(ZX)TY\in\mc I,
	       &
	(XTY)Z & =XT(YZ)\in\mc I.
\end{align}
Thus the obstruction is already restricted to a five-dimensional
algebra with a four-dimensional ideal. The next lemma determines
the active pencil forced by its two orbit dimensions.

\subsection{Structure of the active pencil}

\begin{lemma}
	\label{lem:pencil}
	Let $R>0$, and let $A_0,A_1\in M_2(\C)$ be independent.
	If both spaces in \cref{eq:orbit-matrices} have dimension two,
	there are orthonormal bases $\{\ket{u_0},\ket{u_1}\}$ and
	$\{\ket{v_0},\ket{v_1}\}$ such that
	\begin{align}
		\operatorname{span}\{A_0,A_1\}
		 & =\operatorname{span}
		\{\ket{u_0}\bra{v_0},\ket{u_1}\bra{v_1}\}.
		\label{eq:pencil}
	\end{align}
	Moreover, $R$ is diagonal in both bases, with strictly positive
	diagonal entries.
\end{lemma}

\begin{proof}
	The two-dimensional space
	\begin{align}
		\mc B=\operatorname{span}
		\{\Id_2,A_k^\dagger A_j:j,k=0,1\}
	\end{align}
	contains the identity and is closed under adjoints.
	It therefore contains a nonscalar Hermitian matrix $H$.
	To justify this choice, take any nonscalar element of $\mc B$;
	at least one of its Hermitian real and imaginary parts is nonscalar.
	Since $\dim\mc B=2$, we have
	$\mc B=\operatorname{span}\{\Id_2,H\}$.
	Choose an orthonormal eigenbasis $\{\ket{v_0},\ket{v_1}\}$ of $H$.
	The eigenvalues are distinct, so every matrix in $\mc B$ is diagonal
	in this basis.

	Consequently, for all $j,k$,
	\begin{align}
		\langle A_kv_0,A_jv_1\rangle
		 & =\bra{v_0}A_k^\dagger A_j\ket{v_1}=0.
		\label{eq:column-orthogonality}
	\end{align}
	The column spaces
	\begin{align}
		\mc L_t=\operatorname{span}\{A_0\ket{v_t},A_1\ket{v_t}\},
		\qquad t=0,1,
	\end{align}
	are therefore orthogonal.

	Neither column space can vanish. Suppose, for example, that
	$\mc L_1=0$. Then
	$A_j=\ket{x_j}\bra{v_0}$, and independence of $A_0,A_1$ makes
	$\ket{x_0},\ket{x_1}$ independent. Since $R>0$,
	\begin{align}
		A_jRA_k^\dagger
		 & =\bra{v_0}R\ket{v_0}\ket{x_j}\bra{x_k}.
	\end{align}
	These four matrices span $M_2(\C)$, contradicting the
	two-dimensional column orbit. The case $\mc L_0=0$ is identical
	with the roles of the two columns exchanged.

	The nonzero orthogonal spaces $\mc L_0,\mc L_1$ in $\C^2$ must
	both be lines. Choose unit vectors $\ket{u_t}$ spanning them.
	There are coefficients $d_{jt}$ such that
	\begin{align}
		A_j
		 & =d_{j0}\ket{u_0}\bra{v_0}
		+d_{j1}\ket{u_1}\bra{v_1}.
		\label{eq:pencil-coefficients}
	\end{align}
	The $2\times2$ coefficient matrix $(d_{jt})$ is invertible because
	the two $A_j$ are independent. Hence the two rank-one matrices
	\begin{align}
		K_t & =\ket{u_t}\bra{v_t},\qquad t=0,1,
	\end{align}
	form another basis of the active pencil.

	An invertible change of active basis preserves the span of
	$\{A_jRA_k^\dagger\}$. In the new basis its four generators are
	\begin{align}
		K_sRK_t^\dagger
		 & =\bra{v_s}R\ket{v_t}\ket{u_s}\bra{u_t},
		\qquad s,t=0,1.
		\label{eq:pencil-orbit}
	\end{align}
	The two diagonal coefficients are strictly positive.
	If the off-diagonal coefficient $\bra{v_0}R\ket{v_1}$ were nonzero,
	its conjugate would also be nonzero. All four matrix units in the
	$\ket{u_t}$ basis would then occur in \cref{eq:pencil-orbit}, giving
	a four-dimensional span. This is excluded. Therefore
	\begin{align}
		\bra{v_0}R\ket{v_1}
		 & =\bra{v_1}R\ket{v_0}=0.
	\end{align}
	The span in \cref{eq:pencil-orbit} is exactly the diagonal matrices
	in the $\ket{u_t}$ basis. Since adjoining $R$ leaves its dimension
	equal to two, $R$ belongs to that diagonal space as well.
	Its diagonal entries in either basis are positive by faithfulness.
\end{proof}

The change from $A_j$ to $K_t$ in this proof is an algebraic change
of basis in their span, not an additional physical gauge. It is
used only to describe the same horizontal coordinate space.

\subsection{Failure of closure on the remaining locus}

\begin{theorem}
	\label{thm:no-go}
	Consider the locally purified tensor
	\cref{eq:purification,eq:letters} with Gram ranks $(1,2)$.
	Assume that its trace matrix is idempotent, its periodic operators
	have trace one, and its one-site coefficient span has dimension five.
	Then
	\begin{align}
		\ker F_1\not\subseteq\ker F_2.
		\label{eq:no-kernel-inclusion}
	\end{align}
	In particular, the one-site and two-site coefficient kernels cannot
	be equal. No complex-linear refinement map on the letter span exists,
	and no full-space CPTP maps can satisfy \cref{eq:rfp}.
\end{theorem}

\begin{proof}
	By \cref{prop:faithful}, both transfer weights are faithful.
	Apply the virtual gauge reduction to obtain $\mc E(Y)=R\tr Y$.
	All hypotheses and the kernel inclusion question are invariant
	under this reduction.

	Suppose, for contradiction, that
	$\ker F_1\subseteq\ker F_2$.
	By \cref{eq:kernel-criterion}, $\mc W$ is closed under multiplication.
	By \cref{lem:orbit}, both orbit dimensions equal two.
	Apply \cref{lem:pencil} and use
	$K_t=\ket{u_t}\bra{v_t}$ as a basis of the active pencil.
	Write
	\begin{align}
		R                       & =r_0^u\ket{u_0}\bra{u_0} +r_1^u\ket{u_1}\bra{u_1}
		\nonumber                                                                    \\
		                        & =r_0^v\ket{v_0}\bra{v_0} +r_1^v\ket{v_1}\bra{v_1},
		                        &
		r_0^u,r_1^u,r_0^v,r_1^v & >0.
		\label{eq:two-diagonalizations}
	\end{align}

	Represent a superoperator as a matrix with output basis
	$\ket{u_s}\otimes\overline{\ket{u_t}}$ and input basis
	$\ket{v_s}\otimes\overline{\ket{v_t}}$, both ordered as
	$00,01,10,11$. Denote its matrix in these two bases by a hat.
	This is an invertible change of coordinates for matrix-space
	membership; multiplication of hatted matrices is not assumed to
	represent composition.

	The four active products are
	\begin{align}
		K_s\otimes\overline{K_t}
		 & =
		\left(\ket{u_s}\otimes\overline{\ket{u_t}}\right)
		\left(\bra{v_s}\otimes\overline{\bra{v_t}}\right).
	\end{align}
	Their hatted matrices are the four diagonal matrix units.
	The hatted matrix of $T:Y\mapsto R\tr Y$ is
	\begin{align}
		\widehat T
		 & =
		\begin{pmatrix}
			r_0^u & 0 & 0 & r_0^u \\
			0     & 0 & 0 & 0     \\
			0     & 0 & 0 & 0     \\
			r_1^u & 0 & 0 & r_1^u
		\end{pmatrix}.
		\label{eq:T-mixed-bases}
	\end{align}
	Therefore
	\begin{align}
		\widehat{\mc W}
		 & =\{\Lambda+c\widehat T:
		\Lambda\ \text{diagonal},\ c\in\C\}.
		\label{eq:W-mixed-bases}
	\end{align}

	The matrix $E=K_0\otimes\overline{K_0}$ belongs to $\mc W$.
	Compute its composition with $T$ as a map before changing bases:
	\begin{align}
		(ET)\operatorname{vec}(Y)
		 & =\operatorname{vec}(K_0RK_0^\dagger)\tr Y
		\nonumber                                             \\
		 & =r_0^v\operatorname{vec}(\ket{u_0}\bra{u_0})\tr Y.
	\end{align}
	Thus
	\begin{align}
		\widehat{ET}
		 & =
		\begin{pmatrix}
			r_0^v & 0 & 0 & r_0^v \\
			0     & 0 & 0 & 0     \\
			0     & 0 & 0 & 0     \\
			0     & 0 & 0 & 0
		\end{pmatrix}.
		\label{eq:ET-mixed-bases}
	\end{align}
	If $ET$ belonged to $\mc W$, then
	$\widehat{ET}=\Lambda+c\widehat T$ for some diagonal $\Lambda$.
	The off-diagonal entry with row $11$ and column $00$ would give
	\begin{align}
		0 & =c r_1^u,
		  &
		c & =0.
	\end{align}
	The off-diagonal entry with row $00$ and column $11$ would then give
	\begin{align}
		r_0^v & =c r_0^u=0,
	\end{align}
	contrary to $r_0^v>0$. Hence $ET\notin\mc W$, contradicting closure.

	It follows that $\mc W^2\not\subseteq\mc W$, which proves
	\cref{eq:no-kernel-inclusion}. Any linear refinement map would send
	$F_1(c)=0$ to $F_2(c)=0$, so such a map cannot exist.
	CPTP refinement maps are a special case.
\end{proof}

The proof distinguishes two exhaustive possibilities on the faithful
parameter space. If either transfer orbit has dimension greater
than two, the dimension bound in \cref{lem:orbit} excludes closure.
If both have dimension two, \cref{lem:pencil} applies and the explicit
product \cref{eq:ET-mixed-bases} excludes closure. There is no remaining
coefficient-compatible locus within the stated domain.

\subsection{Increasing only the scalar multiplicity does not help}

\begin{corollary}
	\label{cor:scalar-extension}
	Keep the physical encoding, bond-two local purification, coefficient span
	five, and rank-one idempotent trace transfer. If the active sector Gram
	operator has rank one, no linear refinement exists, irrespective of the
	rank of the scalar sector Gram operator.
\end{corollary}
\begin{proof}
	Write an arbitrary scalar Gram factorization with matrices
	$C_1,\ldots,C_s\in M_2(\C)$. The active sector still has two matrices
	$A_0,A_1$, and the coordinate space remains
	\begin{align}
		\mc W & =\operatorname{span}\{A_i\otimes\overline{A_j},T:i,j=0,1\}, \\
		T     & =\sum_i A_i\otimes\overline{A_i}
		+\sum_{l=1}^{s}C_l\otimes\overline{C_l}.
	\end{align}
	The scalar contributions occur through one physical coordinate, not
	through $s$ independent physical coordinates. If either factor in
	$\mc E(X)=R\tr(LX)$ is singular, positivity of its Choi matrix forces
	all Kraus matrices into a common input or output line. Their coordinate
	span then has dimension at most four, a contradiction.
	Otherwise the same virtual gauge gives $R>0$, $L=I$, and $\tr R=1$.
	Full span five implies independence of $A_0,A_1$.
	The orbit and pencil arguments and the final product contradiction in
	\cref{thm:no-go} use only these properties and the displayed coordinate
	space. They do not use the number of scalar Kraus matrices.
	Thus $\mc W^2\not\subseteq\mc W$, excluding refinement.
\end{proof}

The orthonormal four-Kraus parameterization in \cref{sec:transfer}
is not being extended to nonminimal Kraus families. Only the obstruction
extends. In this encoding the scalar Gram operator acts on $\C^4$, so
its rank is at most four. Consequently none of the pairs $(1,1)$,
$(1,2)$, $(1,3)$, or $(1,4)$ can pass the trace-one, full-span test.
Since the purification rank obeys
\begin{align}
	r_E\geq2\rank B_0+\rank B_1,
\end{align}
the first unexcluded pair is $(2,1)$, requiring $r_E\geq5$.
This identifies a necessary change in the ansatz, not a feasible example.

\subsection{Pauli regression and physical channel requirements}

The accepted regression seed is
\begin{align}
	A_0=\Id_2/2,\qquad A_1=X/2,\qquad
	C_0=Y/2,\qquad C_1=Z/2,
	\label{eq:pauli}
\end{align}
where $X,Y,Z$ are the Pauli matrices.
It has $R=\Id_2/2$ and active Hilbert--Schmidt plane
$\operatorname{span}\{\Id_2,X\}$.
Its one-site coordinates span
\begin{align}
	\Id_2\otimes\Id_2,\quad
	\Id_2\otimes X,\quad
	X\otimes\Id_2,\quad
	X\otimes X,\quad
	Q_0:=Y\otimes\overline Y+Z\otimes\overline Z.
	\label{eq:pauli-span}
\end{align}
An active--scalar product produces
\begin{align}
	P_{01}Q
	 & =\sum_{l=0}^1A_0C_l\otimes\overline{A_1C_l}
	\nonumber                                      \\
	 & =\frac1{16}\left(Y\otimes\overline{iZ}
	+Z\otimes\overline{-iY}\right)
	\nonumber                                      \\
	 & =-\frac{i}{16}(Y\otimes Z+Z\otimes Y).
	\label{eq:pauli-product}
\end{align}
The identities used here are $XY=iZ$, $XZ=-iY$,
$\overline Y=-Y$, and $\overline Z=Z$.
Pauli tensor orthogonality makes this nonzero coordinate independent
of \cref{eq:pauli-span}. Thus closure fails directly.
The exact regression script~\cite{Checks} verifies the stronger
finite statement that the full blocked span is six-dimensional,
whereas the one-site span is five-dimensional.
The general theorem does not assume a numerical search result or
infer a universal blocked rank from this example.

The obstruction in \cref{thm:no-go} occurs before complete positivity.
In general, equality of coefficient kernels would provide a linear
isomorphism between letter spans, not a quantum channel on the full
physical spaces. Such channels would require positive Choi matrices
$J_{\mc T}$ and $J_{\mc S}$ obeying
\begin{align}
	J_{\mc T}                   & \geq0,
	                            &
	\tr_{\mathrm{out}}J_{\mc T} & =\Id_5,
	\nonumber                                \\
	J_{\mc S}                   & \geq0,
	                            &
	\tr_{\mathrm{out}}J_{\mc S} & =\Id_{25},
	\label{eq:channel-constraints}
\end{align}
together with all letter equations in \cref{eq:rfp}.
Both Choi matrices would have size $125\times125$.

The stronger accepted encoding ansatz additionally asks for
normalized faithful states $\sigma_0\in M_r$, $\sigma_1\in M_s$
and an isometry
$V:(\C^2\otimes\C^r)\oplus\C^s\to\mc H^{\otimes2}$ such that
\begin{align}
	J_2(X\oplus z)
	    & =V[(X\otimes\sigma_0)\oplus z\sigma_1]V^\dagger,
	\nonumber                                              \\
	G^{\alpha\beta}
	    & =J_2(X_{\alpha\beta}),
	    &
	r,s & \geq1,\qquad 2r+s\leq25.
	\label{eq:stationarity}
\end{align}
Compression and preparation can provide full-space decoders for
such encodings~\cite{Kato}, but they cannot establish the second
line of \cref{eq:stationarity}. In the present stratum that line
already fails as a linear assignment, by \cref{thm:no-go}.
No choice of $r,s,V,\sigma_0,\sigma_1$ repairs the coefficient
obstruction.
